\section{離散双対は軌道積分ではない}
\label{sec:discrete}

地位は\dstrec{}である。ソースは \texttt{Algebra.Discrete} である。この章の手続きは時間ループに入らない。

\subsection{格子}

正整数 \(N\) に対し、ねじれの標本は \((\mathbb{Z}/N\mathbb{Z})^6\) である。
\begin{verbatim}
struct DiscreteTorsion:
    n: array[3] of mod N   # usual
    m: array[3] of mod N   # dual
\end{verbatim}
点数は \(N^6\) である。連続パラメータへの埋め込みは
\[
\alpha_a=\frac{2\pi}{N}(n_a\bmod N),\qquad
\beta_a=\frac{2\pi}{N}(m_a\bmod N)
\]
であり、識別子は \texttt{toTorsionParams} である。値は \([0,2\pi)\) に入る。許容かどうかは、埋め込んだあとの \texttt{IsAdmissibleContinuous} で判定する。格子点であること自体は許容を意味しない。

\subsection{トーラスは模型選択である}

循環指数は \(2\pi\) 周期である。双曲指数は \(2\pi\) 周期ではない。ブーストをこのトーラスへ載せることは、離散化の模型選択である。許容錐は、トーラスに載せたことから自動では従わない。列挙コマンドは、選択であることをヘルプ文字列に書く。

\subsection{コマンド}

\begin{verbatim}
command enumerate_admissible(N) -> list of DiscreteTorsion:
    for each of the N^6 points:
        p = to_torsion_params(point)
        if admissible(p): emit point
\end{verbatim}

出力は探索の入力である。各点の \(J\) と \(M\) を添えてよい。点を時刻順に並べ、速度として \texttt{working\_step} へ渡さない。\(N\) を段数と解釈しない。

埋め込みの成分は常に非負なので、離散点に限れば \texttt{Discrete.IsAdmissible} と、埋め込んだあとの \texttt{IsPrincipalBranch} は同値である。連続状態では同値ではない。列挙の検査は \texttt{IsAdmissible} を呼ぶ。

双対角だけを \((2\pi/N)\mathbb{Z}\) に制限し、通常ラピディティを実数のまま残す読みは、上の六次元格子とは別の模型である。同じコマンドにまとめない。
